% ===========================================================================
% CHAPTER 2
% ===========================================================================

\chapter{Technical Background and Related Work}
\label{chap:background-related-work}

\section{TLS and the X.509 Certificate Ecosystem}
\label{sec:tls-x509-ecosystem}

\textbf{Transport Layer Security (TLS)} is the standard protocol for protecting
Internet communication through \textbf{confidentiality}, \textbf{integrity},
and \textbf{authentication}. \textbf{HTTPS} applies TLS to web traffic. During
the handshake, the client and server negotiate cryptographic parameters,
establish shared secrets, authenticate the server, and derive the symmetric session keys used
to protect subsequent application data~\cite{rfc8446}.

\begin{figure}[H]
  \centering
  \begin{tikzpicture}[
    font=\scriptsize,
    >=Latex,
    participant/.style={draw, rounded corners=2pt, align=center, minimum width=2.4cm, minimum height=0.7cm, fill=black!5},
    life/.style={densely dashed, black!55},
    message/.style={-{Latex[length=2mm]}, line width=0.55pt},
    secure/.style={{Latex[length=2mm]}-{Latex[length=2mm]}, line width=0.65pt},
    messageLabel/.style={fill=white, inner sep=1.3pt, align=center},
    annotation/.style={fill=white, inner sep=1.5pt, align=center, text width=3.8cm}
  ]
    \node[participant] (client) at (-4.8,0) {\textbf{Client}};
    \node[participant] (server) at (4.8,0) {\textbf{Server}};

    \draw[life] (client.south) -- ++(0,-5.1);
    \draw[life] (server.south) -- ++(0,-5.1);

    \draw[message] (-4.8,-0.9) --
      node[messageLabel,above]{\textbf{ClientHello}}
      (4.8,-0.9);

    \draw[message] (4.8,-1.85) --
      node[messageLabel,above]{\textbf{ServerHello}, \textbf{Certificate},\\
      \textbf{CertificateVerify}, \textbf{Finished}}
      (-4.8,-1.85);

    \node[annotation] (validation) at (-3.25,-2.75)
      {\textbf{Client validates certificate}\\and trust chain};
    \node[annotation] (keys) at (0,-3.85)
      {\textbf{Shared session keys established}};

    \draw[message] (-4.8,-3.2) -- (keys.west);
    \draw[message] (4.8,-2.45) -- (keys.east);

    \coordinate (httpsPoint) at (0,-5.05);
    \draw[message] (keys.south) -- (httpsPoint);
    \draw[secure] (-4.8,-5.05) -- (4.8,-5.05);
    \node[messageLabel,font=\scriptsize\bfseries] at (httpsPoint)
      {Encrypted HTTPS communication};
  \end{tikzpicture}
  \caption{Simplified TLS Handshake and Certificate Authentication Process}
  \label{fig:tls-handshake-authentication}
\end{figure}

\textbf{Public Key Infrastructure (PKI)} provides the trust framework used for
certificate-based authentication. A \textbf{root Certificate Authority (CA)}
acts as a trust anchor and delegates issuance through one or more
\textbf{intermediate CAs}. Together with the server certificate, these
certificates form a \textbf{trust chain}. A browser validates the signatures,
validity periods, identity constraints, and trusted root before accepting the
server's identity.

The \textbf{X.509} standard defines the structure of public-key certificates.
Important fields include the \textbf{Serial Number}, \textbf{Subject},
\textbf{Issuer}, \textbf{Validity Period}, \textbf{Signature Algorithm}, and
\textbf{Subject Public Key Info}. Extensions such as
\textbf{Basic Constraints}, \textbf{Key Usage}, \textbf{Extended Key Usage},
\textbf{Certificate Policies}, and \textbf{Subject Alternative Name (SAN)}
describe certification authority status, permitted operations, policy context,
and authenticated domain names~\cite{rfc5280}. Collectively, these fields define
a certificate's identity, cryptographic properties, and intended use.

Beyond their role in authentication, X.509 certificates provide a rich source 
of publicly observable security information. Their fields support the study of 
cryptographic practices, issuer behaviour, deployment patterns, and certificate 
quality. Comparing these properties across large datasets enables ecosystem-wide measurements and
Internet security research that cannot be achieved through isolated
certificate inspection.

\textbf{Certificate observation} differs from \textbf{certificate validation}.
Browsers trust a connection only after validation, while research measurements
may record expired, self-signed, hostname-mismatched, or untrusted
certificates. Observation confirms presentation by an endpoint, not acceptance
by a standards-compliant client.

\section{Certificate Transparency and Dataset Freshness}
\label{sec:ct-freshness}

\textbf{Certificate Transparency (CT)} was introduced to reduce the risk that
incorrectly or maliciously issued certificates could remain undetected within
the \textbf{Public Key Infrastructure (PKI)}. It requires certificate issuance
to be exposed through publicly accessible, cryptographically verifiable, and
\textbf{append-only logs}. This public record improves the transparency and
accountability of Certificate Authorities by allowing certificate activity and
log consistency to be independently examined~\cite{rfc9162}.

Within the CT ecosystem, a \textbf{Certificate Authority (CA)} submits an
issued certificate or precertificate to a CT log. The log records the entry and
returns a \textbf{Signed Certificate Timestamp (SCT)} as evidence of promised
inclusion. Browsers check for acceptable SCT evidence when evaluating a
certificate, while \textbf{monitors} and \textbf{auditors} observe the public
logs for unexpected certificates and inconsistent log behaviour. CT therefore
adds public accountability to certificate issuance without replacing normal
certificate-chain and hostname validation.

\begin{figure}[H]
  \centering
  \begin{tikzpicture}[
    font=\scriptsize,
    >=Latex,
    entity/.style={draw, rounded corners=2pt, align=center, text width=3.15cm, minimum height=0.85cm, fill=black!5},
    primary/.style={draw, very thick, rounded corners=2pt, align=center, text width=3.35cm, minimum height=0.9cm, fill=black!9},
    result/.style={draw, rounded corners=2pt, align=center, text width=3.4cm, minimum height=0.85cm, fill=black!4},
    flow/.style={-{Latex[length=2mm]}, line width=0.55pt}
  ]
    \node[entity] (ca) at (0,0) {\textbf{Certificate Authority}\\issues and submits certificate};
    \node[primary] (log) at (0,-1.75) {\textbf{Certificate Transparency Log}\\append-only public record};

    \node[entity] (browser) at (-4.5,-3.65) {\textbf{Browser}\\checks SCT evidence};
    \node[entity] (monitor) at (0,-3.65) {\textbf{Monitor or Auditor}\\observes log activity};
    \node[entity] (researcher) at (4.5,-3.65) {\textbf{Security Researcher}\\discovers certificate events};

    \node[result] (analysis) at (0,-5.35) {\textbf{Certificate Analysis}\\monitoring, discovery, and measurement};

    \draw[flow] (ca) -- (log);
    \draw[flow] (log) -- (browser);
    \draw[flow] (log) -- (monitor);
    \draw[flow] (log) -- (researcher);
    \draw[flow] (monitor) -- (analysis);
    \draw[flow] (researcher) -- (analysis);
  \end{tikzpicture}
  \caption{Simplified Certificate Transparency Ecosystem}
  \label{fig:certificate-transparency-ecosystem}
\end{figure}

The public visibility provided by CT supports \textbf{certificate monitoring},
\textbf{certificate discovery}, and the detection of potential
\textbf{misissuance}. Organizations can search for unexpected certificates
associated with their domains, while researchers can examine issuance
patterns, CA behaviour, cryptographic characteristics, and changes across large
collections. CT logs consequently provide an important data source for
Internet-scale measurement and certificate ecosystem research.

\textbf{Dataset freshness} describes how closely a dataset represents the
current state of the environment being studied. Certificate ecosystems change
continuously as new certificates are issued, existing certificates expire or
are renewed, new domains become active while others are retired or no longer deploy TLS services, and deployed certificates are
replaced. A static certificate dataset therefore becomes progressively
outdated and may no longer represent current issuance or deployment patterns.

Maintaining a fresh dataset is essential for accurate large-scale analytics,
meaningful statistical comparisons, reliable trend analysis, longitudinal
security studies, and informed security assessment. Historical snapshots remain
valuable for studying change, but current measurements require continuing
observation of the evolving certificate ecosystem. Continuously updated
datasets are therefore fundamental to meaningful SSL/TLS certificate
analytics.

\section{Security-Relevant Certificate Characteristics}
\label{sec:certificate-key-reuse}

An \textbf{X.509 certificate} contains substantially more information than a
binding between a domain identity and a public key. Its subject, issuer,
validity, cryptographic parameters, extensions, and policy information describe
how the certificate was issued and how it is intended to be used~\cite{rfc5280}.
These attributes provide observable evidence about \textbf{deployment
practices}, \textbf{cryptographic strength}, \textbf{operational security}, and
broader certificate ecosystem behaviour. They are therefore widely used in
Internet measurement studies and cybersecurity research.

\begin{figure}[H]
  \centering
  \begin{tikzpicture}[
    font=\fontsize{9}{10.8}\selectfont,
    >=Latex,
    root/.style={draw, very thick, rounded corners=2pt, align=center, text width=3.5cm, minimum height=0.95cm, fill=black!9},
    identityCategory/.style={draw, rounded corners=2pt, align=center, text width=4.1cm, minimum height=1.9cm, inner xsep=3pt, fill=black!5},
    cryptographyCategory/.style={draw, rounded corners=2pt, align=center, text width=2.9cm, minimum height=1.9cm, inner xsep=3pt, fill=black!5},
    operationalCategory/.style={draw, rounded corners=2pt, align=center, text width=2.9cm, minimum height=1.9cm, inner xsep=3pt, fill=black!5},
    trustCategory/.style={draw, rounded corners=2pt, align=center, text width=3.2cm, minimum height=1.9cm, inner xsep=3pt, fill=black!5},
    indicator/.style={draw, rounded corners=2pt, align=center, text width=10.5cm, minimum height=1.45cm, fill=black!7},
    outcome/.style={draw, very thick, rounded corners=2pt, align=center, text width=4.5cm, minimum height=0.9cm, fill=black!9},
    flow/.style={-{Latex[length=2mm]}, line width=0.55pt},
    connector/.style={line width=0.5pt}
  ]
    \node[root] (certificate) at (0,0) {\textbf{X.509 Certificate}};

    \node[identityCategory] (identity) at (-5.15,-2.0)
      {\textbf{Identity}\\Subject\\Issuer\\
      \mbox{\fontsize{8.5}{10}\selectfont Subject Alternative Name (SAN)}};
    \node[cryptographyCategory] (crypto) at (-1.22,-2.0)
      {\textbf{Cryptography}\\RSA or ECDSA\\Key length\\
      \mbox{Signature algorithm}};
    \node[operationalCategory] (operational) at (2.12,-2.0)
      {\textbf{Operational}\\Validity period\\Expiration date\\Certificate lifetime};
    \node[trustCategory] (trust) at (5.60,-2.0)
      {\textbf{Trust}\\\mbox{Certificate Authority}\\Certificate chain\\Validation level};

    \draw[connector] (certificate.south) -- (0,-0.85);
    \draw[connector] (-5.15,-0.85) -- (5.60,-0.85);
    \draw[flow] (-5.15,-0.85) -- (identity.north);
    \draw[flow] (-1.22,-0.85) -- (crypto.north);
    \draw[flow] (2.12,-0.85) -- (operational.north);
    \draw[flow] (5.60,-0.85) -- (trust.north);

    \node[indicator] (indicators) at (0,-4.45)
      {\mbox{\textbf{Security-Relevant Indicators}}\\[2pt]
      Weak cryptography \quad \textbullet \quad Expired certificates \quad
      \textbullet \quad Public-key reuse\\
      Certificate quality \quad \textbullet \quad Standards-compliance  };

    \draw[connector] (identity.south) -- (-5.15,-3.35);
    \draw[connector] (crypto.south) -- (-1.22,-3.35);
    \draw[connector] (operational.south) -- (2.12,-3.35);
    \draw[connector] (trust.south) -- (5.60,-3.35);
    \draw[connector] (-5.15,-3.35) -- (5.60,-3.35);
    \draw[flow] (0,-3.35) -- (indicators.north);

    \node[outcome] (assessment) at (0,-6.2)
      {\mbox{\textbf{Security Assessment}}};
    \draw[flow] (indicators) -- (assessment);
  \end{tikzpicture}
  \caption{Major Security-Relevant Characteristics of an X.509 Certificate}
  \label{fig:x509-security-characteristics}
\end{figure}

\textbf{Cryptographic characteristics} determine the mechanisms used to
authenticate the certificate and protect the associated connection. The
\textbf{public-key algorithm}, commonly RSA or ECDSA, identifies the underlying
cryptographic approach, while \textbf{key length} affects resistance to
computational attacks. The \textbf{signature algorithm} identifies how the
issuer signs the certificate and which hash function supports that signature.
Modern algorithms and appropriately sized keys contribute to strong security,
whereas obsolete hashes, weak algorithms, or insufficient key sizes may reduce
confidence in a deployment. These characteristics must also be interpreted in
historical and policy context, because an algorithm accepted when a certificate
was issued may later become deprecated as standards and computational
capabilities evolve. Algorithm choice and key size therefore indicate whether
a deployment follows contemporary security expectations.

\textbf{Operational characteristics} describe certificate lifecycle and trust
management. The \textbf{validity period} defines when a certificate is intended
to be accepted, and the expiration date indicates when it should cease to be
used. Expired certificates can interrupt services or reveal weaknesses in
renewal and deployment practices, while unusually long validity periods may
increase exposure if authorization or key ownership changes~\cite{ma2023stale}.
The issuing \textbf{Certificate Authority} and certificate chain also provide
context about the trust relationship, validation process, and certificate
management practices. Issuer patterns across large datasets can additionally
reveal trust concentration, differing validation practices, and changes in CA
market activity.

Common \textbf{security-relevant indicators} include expired certificates,
weak cryptographic configurations, repeated public keys, and
standards-compliance or certificate-quality findings. Public-key reuse can
arise across distinct certificates and has been observed at Internet
scale~\cite{nezhadian2025reuse}; however, legitimate renewal, cross-signing,
and shared infrastructure can produce the same pattern
~\cite{junaid2025shared}. These characteristics should therefore be treated as
signals requiring contextual investigation, not as automatic proof of a
vulnerability, successful attack, or private-key compromise. When considered
together, identity, cryptographic, operational, and trust attributes provide a
broader view of SSL/TLS security posture and support meaningful ecosystem-wide
measurement.

\section{Related Work and Research Gap}
\label{sec:related-work-comparison}

SSL/TLS certificates, \textbf{Certificate Transparency (CT)}, and the broader
Web Public Key Infrastructure have been studied extensively through active
scanning, passive observation, public log analysis, and longitudinal
measurement. Previous work generally falls into three overlapping categories:
\textbf{Internet-scale certificate measurement}, \textbf{CT monitoring}, and
\textbf{certificate search and analysis platforms}. Collectively, these
efforts have improved visibility into certificate deployment, trust
relationships, issuance quality, cryptographic practices, and changes in the
certificate ecosystem. Table~\ref{tab:related-work-comparison} summarizes
representative studies most relevant to large-scale certificate analytics.

\begin{table}[H]
  \centering
  \caption{Representative Studies on SSL/TLS Certificate Analysis}
  \label{tab:related-work-comparison}
  \begingroup
  \scriptsize
  \renewcommand{\arraystretch}{1.12}
  \setlength{\tabcolsep}{3pt}
  \begin{tabularx}{\textwidth}{@{}>{\RaggedRight\arraybackslash}p{0.17\textwidth}>{\RaggedRight\arraybackslash}p{0.18\textwidth}>{\RaggedRight\arraybackslash}p{0.27\textwidth}>{\RaggedRight\arraybackslash}X@{}}
    \toprule
    \textbf{Study} & \textbf{Primary Focus} & \textbf{Major Contribution} & \textbf{Relevance to This Work} \\
    \midrule
    Holz et al.~\cite{holz2011ssl}
    & Active and passive measurement
    & Characterized X.509 deployment and trust using complementary observation methods.
    & Establishes the value of ecosystem-wide certificate measurement. \\

    Durumeric et al.~\cite{durumeric2013https}
    & Longitudinal HTTPS measurement
    & Used repeated Internet-wide scans to examine CA relationships, deployment practices, and configuration problems.
    & Motivates large-scale and time-aware certificate analytics. \\

    Laurie et al.~\cite{rfc9162}
    & Certificate Transparency
    & Formalized publicly auditable logs for observing certificate issuance and detecting suspect activity.
    & Provides the foundation for CT-based discovery and dataset freshness. \\

    VanderSloot et al.~\cite{vandersloot2016complete}
    & Measurement coverage
    & Compared CT logs, active scans, domain crawls, and passive observations, showing that the sources provide complementary views.
    & Supports combining public observations with endpoint measurement. \\

    Kumar et al.~\cite{kumar2018misissuance}
    & Certificate quality and CA behaviour
    & Introduced ZLint and measured certificate conformance with technical and policy requirements.
    & Supports certificate-quality indicators and comparative issuer analysis. \\

    Ma et al.~\cite{ma2023stale}
    & Certificate lifecycle
    & Examined valid but operationally stale certificates and the security implications of prolonged key validity.
    & Motivates validity, lifetime, and freshness analysis. \\

    Nezhadian et al.~\cite{nezhadian2025reuse}
    & Public-key reuse
    & Measured the propagation of reused RSA certificates and keys across large PKI datasets.
    & Supports fingerprint-based analysis of shared public keys. \\
    \bottomrule
  \end{tabularx}
  \endgroup
\end{table}

Public platforms complement academic measurement by making Internet and
certificate observations searchable. \textbf{crt.sh} primarily supports
certificate search and CT-log exploration~\cite{crtsh}; \textbf{Censys}
combines Internet-wide observations with searchable certificate records,
derived fields, and statistical reporting~\cite{durumeric2015censys,censyscerts}; and
\textbf{Shodan} indexes Internet-connected services and exposes TLS
certificate fields through search, filters, facets, reports, and
maps~\cite{shodan,shodanct}. These services are valuable for discovery, monitoring, and
targeted investigation, although their analytical emphasis and presentation
differ. Table~\ref{tab:platform-comparison} therefore uses descriptive
capability levels rather than treating the platforms as directly equivalent.

\begin{table}[H]
  \centering
  \caption{Comparison of Representative SSL/TLS Certificate Platforms}
  \label{tab:platform-comparison}
  \begingroup
  \scriptsize
  \renewcommand{\arraystretch}{1.14}
  \setlength{\tabcolsep}{2.5pt}
  \begin{tabularx}{\textwidth}{@{}>{\RaggedRight\arraybackslash}p{0.14\textwidth}>{\RaggedRight\arraybackslash}X>{\RaggedRight\arraybackslash}X>{\RaggedRight\arraybackslash}X>{\RaggedRight\arraybackslash}X>{\RaggedRight\arraybackslash}X@{}}
    \toprule
    \textbf{Platform} & \textbf{Certificate Search} & \textbf{CT Support} & \textbf{Security Analytics} & \textbf{Continuous Dataset Updates} & \textbf{Interactive Visualization} \\
    \midrule
    crt.sh
    & Certificate and domain queries
    & Primary focus
    & Certificate details and lint findings
    & Continuous CT-log monitoring
    & Search and record views \\

    Censys
    & Structured certificate search
    & CT and scan-derived records
    & Trust, validation, lint, and host context
    & Ongoing CT ingestion and Internet scans
    & Search and statistical reports \\

    Shodan
    & TLS fields within service search
    & Separate CT lookup support
    & TLS filters, facets, and selected exposure indicators
    & Recurring Internet-wide observations
    & Reports, maps, and search views \\

    SSL Guardian
    & Certificate search, filtering, and detail views
    & CT-assisted discovery
    & Certificate indicators, shared-key analysis, and CA comparison
    & Cycle-based acquisition and analytical refresh
    & Certificate-focused dashboards, charts, tables, and drill-down \\
    \bottomrule
  \end{tabularx}
  \endgroup
\end{table}

Existing studies commonly address a specific measurement question, data
source, security property, or issuer behaviour, while public platforms
primarily emphasize certificate lookup, Internet asset search, CT visibility,
or general-purpose monitoring. Several platforms already combine substantial
parts of this functionality; therefore, the remaining gap is not the absence
of certificate-search or Internet-measurement tools. Rather, relatively
limited attention has been given to a certificate-focused analytical
environment that brings together \textbf{automated certificate acquisition},
\textbf{CT-assisted dataset maintenance}, \textbf{large-scale descriptive
analytics}, \textbf{security-oriented indicators}, \textbf{Certificate
Authority comparison}, and \textbf{interactive visualization} within one
coherent workflow. This systems-integration gap provides the motivation for
SSL Guardian while preserving an important distinction: its analytical
indicators support comparative exploration and further investigation, rather
than replacing formal certificate validation, security auditing, or
externally calibrated risk assessment.
